% PDF p. 301
\section{频率混沌游戏表示（FCGR）矩阵与多位点遗传关联研究}

我们将多位点基因型序列转换成 FCGR 矩阵。然后用两组的矩阵构造一个 $2\times N$ 列联表，用统计检验考察与疾病的关联或列联表单元格计数的独立性。这里还要问，能否用格拉斯曼流形来考察这些数据？下面讨论这种方法。

% PDF p. 302
简单来说，格拉斯曼流形描述的是 $\R^n$ 的 $k$ 维子空间。在 CGR 的语境中，可以进行降维，也可以把 FCGR 矩阵各列张成的子空间视为格拉斯曼流形上的一个点。此外，FCGR 矩阵本身表示基因型频率，因而可以转换为列联表，用于对疾病关联进行统计检验。

下面说明如何将FCGR 矩阵得到的双向列联表推广到格拉斯曼流形，以及如何在这一背景下检验关联或独立性。

\subsection{从基因型到 FCGR 和列联表}

假设从患病与健康两个群体的多位点基因型序列出发。混沌游戏表示将序列转换成类似分形的图像，其频率变体 FCGR 则以矩阵形式记录 $k$-mer 或其他特征的相对频率。对每个群体，都能得到一个汇总基因型频率模式的 FCGR 矩阵。

现在，将这些 FCGR 矩阵按对应位置排列，即可把它们“展平”或汇总成一个 $2\times n$ 列联表。

\textbf{行}\quad 对应患病与健康两类。

\textbf{列}\quad 表示 FCGR 中的 $n$ 个单元格或特征，可对应特定的 $k$-mer 模式、区域或频率分组。

每个单元格 $(i,j)$ 的计数等于特征 $j$ 在群体 $i$ 中出现的次数。

\subsection{用格拉斯曼流形描述数据}

一个 $2\times n$ 列联表 $C$（$i=1,2$，$j=1,\ldots,n$）可视为一个矩阵，其两行分别表示两个群体的计数分布。如果这两行线性无关，$C$ 的行空间就是 $\R^n$ 的二维子空间。换言之，该表在格拉斯曼流形（见第 1 章）
\[
  \Gr(2,n)
\]
上定义了一个点。这个流形由 $\R^n$ 的所有二维线性子空间组成。

如果希望在考虑子空间几何的情况下比较群体间频率分布的整体“形状”或结构，这一视角尤为有用。它还允许我们利用格拉斯曼流形上的微分几何工具，例如测地距离和基于流形的统计量，检验群体之间的差异。

\subsection{用 \texorpdfstring{$2\times n$}{2×n} 列联表统计量检验关联}

我们考虑两种方法，即欧氏空间中经典的 $2\times n$ 列联表检验和格拉斯曼流形上的测地距离检验。

\subsubsection{经典方法}

用 $2\times n$ 列联表检验基因与疾病关联的一种标准方法是皮尔逊卡方检验。若 $n_{ij}$ 表示单元格 $(i,j)$ 的观测计数，且在独立性（无关联）零假设下，期望计数为
% PDF p. 303
\[
  E_{ij}=\frac{n_{i\cdot}n_{\cdot j}}{n},
\]
则卡方统计量为
\[
  X^2=\sum_{i=1}^2\sum_{j=1}^n\frac{(n_{ij}-E_{ij})^2}{E_{ij}},
\]
在标准正则性条件下，该统计量渐近服从自由度为 $(2-1)(n-1)=n-1$ 的卡方分布。

\subsubsection{基于格拉斯曼流形的方法}

如前所述，将 $2\times n$ 列联表视为 $\Gr(2,n)$ 上的点后，可以采用几何视角。例如，假设在独立性（或无关联）的零假设下，两个群体的计数分布在子空间表示或特征方面“相似”。于是可考虑一种检验统计量，它基于格拉斯曼流形上观测子空间 $S_{\mathrm{obs}}$ 与独立或无关联模型所预测的“零假设”子空间 $S_0$ 之间的测地距离。

常用步骤如下。
\begin{enumerate}[label=步骤 \arabic*.,leftmargin=*]
  \item \textbf{估计零假设子空间}\par
  在无关联的零假设下，可以假定两个 FCGR 矩阵（或相应的频率分布）具有相同的分布。合并后的列联表或相应的估计在 $\R^n$ 中定义了一个“零假设”二维子空间 $S_0$。\translatornote{若零假设下直接令两组的频率行成比例，则相应的行空间只有一维。此时还需定义如何得到用于距离比较的二维参照子空间。}

  \item \textbf{计算观测子空间}\par
  观测列联表的两行定义了观测到的二维子空间 $S_{\mathrm{obs}}$。

  \item \textbf{度量差异}\par
  格拉斯曼流形上的测地距离 $d_G(S_{\mathrm{obs}},S_0)$ 可用两个子空间之间的主角计算。若主角为 $\theta_1$ 和 $\theta_2$，一个常用距离为（附录 6I）
  \[
    d_G(S_{\mathrm{obs}},S_0)=\sqrt{\theta_1^2+\theta_2^2}.
  \]

  \item \textbf{检验统计量及其渐近分布}\par
  零假设下，可借助随机矩阵理论和格拉斯曼流形的几何，推导 $d_G(S_{\mathrm{obs}},S_0)$ 或它的某个函数的渐近分布。可以使用置换检验，即随机重新分配行标签，也可以推导近似的渐近分布，从而获得 $p$ 值。

  例如，若适当缩放的距离平方的渐近分布为卡方随机变量的加权和（附录 6I），
  \[
    nd_G^2(S_{\mathrm{obs}},S_0)\xrightarrow{d}\sum_{i=1}^2\lambda_i\chi_{1,i}^2,
  \]
% PDF p. 304
  且权重 $\lambda_i$ 取决于 $\Gr(2,n)$ 的曲率和底层变异性，就可以将观测统计量与这一参照分布比较，以计算 $p$ 值。
\end{enumerate}

这个框架将经典分类数据分析与现代几何数据分析联系起来，也使我们能够用微分几何工具，例如格拉斯曼流形上的测地距离，检验高维生物数据中的关联。

\paragraph{步骤小结}

\begin{enumerate}
  \item 计算 $U^{\mathsf T}V=W\Sigma Z^{\mathsf T}$。
  \item 取出 $\sigma_i$，令 $\theta_i=\arccos\sigma_i$。
  \item 距离为 $d_G=\lVert\theta\rVert_2$。
  \item 统计量为 $T^2=nd_G^2$。
  \item 渐近关系为 $\sqrt{\theta}\approx N(0,\Sigma_\theta)\Rightarrow T^2\to\sum\lambda_r\chi_1^2$。\translatornote{原印式的根号内为 $\theta$。按附录 6I 的缩放及上一项的统计量定义，此处使用 $\sqrt n\,\theta$。}
\end{enumerate}
至此，得到了主角、格拉斯曼测地距离、检验统计量及其渐近加权卡方分布。

\subsection{第一类错误率}

为确保检验有效，需要计算第一类错误率（type I error rate）。我们考虑三种检验统计量，即欧氏 $\chi^2$、流形加权 $\chi^2$ 和格拉斯曼流形（置换）。取三个名义水平 $\alpha=0.05,0.01,10^{-5}$，即在全局零假设下，每种检验都有 $\Pr(\text{拒绝 }H_0)=\alpha$，与样本量无关。表 \ref{tab:6.4} 给出三种检验统计量的第一类错误率。

\begin{table}[htbp]
  \centering\small
  \TableCaption{6.4}{三种检验统计量的第一类错误率}
  \begin{tabular}{ccccc}
    \toprule
    样本量 & 名义 $\alpha$ & 欧氏 $\chi^2$ & 流形加权 $\chi^2$ & 格拉斯曼流形（置换）\\
    \midrule
    500  & 0.05 & 0.05 & 0.05 & 0.05\\
    500  & 0.01 & 0.01 & 0.01 & 0.01\\
    500  & $1\times10^{-5}$ & $1\times10^{-5}$ & $1\times10^{-5}$ & $1\times10^{-5}$\\
    1000 & 0.05 & 0.05 & 0.05 & 0.05\\
    1000 & 0.01 & 0.01 & 0.01 & 0.01\\
    1000 & $1\times10^{-5}$ & $1\times10^{-5}$ & $1\times10^{-5}$ & $1\times10^{-5}$\\
    2000 & 0.05 & 0.05 & 0.05 & 0.05\\
    2000 & 0.01 & 0.01 & 0.01 & 0.01\\
    2000 & $1\times10^{-5}$ & $1\times10^{-5}$ & $1\times10^{-5}$ & $1\times10^{-5}$\\
    5000 & 0.05 & 0.05 & 0.05 & 0.05\\
    5000 & 0.01 & 0.01 & 0.01 & 0.01\\
    5000 & $1\times10^{-5}$ & $1\times10^{-5}$ & $1\times10^{-5}$ & $1\times10^{-5}$\\
    \bottomrule
  \end{tabular}
\end{table}

\textbf{为什么这些水平能保持精确}
\begin{enumerate}
  \item \textbf{欧氏 $\chi^2$} 按定义，零假设下皮尔逊统计量收敛到 $\chi^2$，因而在极限中 $\Pr(\chi_\nu^2>\chi_{\mathrm{obs}}^2)=\alpha$。即使在 $n\geq 500$ 时，两者也很接近。
% PDF p. 305
  \item \textbf{流形加权 $\chi^2$} 我们推导出精确的混合卡方零分布，使得 $\Pr(\sum_r\lambda_r\chi_{1r}^2>\chi_{\mathrm{obs}}^2)=\alpha$。
  \item \textbf{格拉斯曼流形（置换）} 在零假设下置换标签，得到基于可交换性的精确检验。使用所有随机置换（或足够多次随机置换）时，有 $\Pr(p_{\mathrm{perm}}\leq\alpha)=\alpha$。
\end{enumerate}
总之，三种方法都有效，将第一类错误控制在名义水平 $\alpha=0.05,0.01,10^{-5}$，与每组样本量 $n$ 无关。\translatornote{表 \ref{tab:6.4} 按原值照录。渐近零分布本身不能保证有限样本的拒绝率等于名义水平，通常的非随机化置换检验也可因可取 $p$ 值离散而使拒绝率低于名义水平。}
